\section{A direct-completion obstruction}\label{app:completion}

The following example explains why the contribution of added edges must
be controlled in a completeness-preserving reduction.

\begin{proposition}
Completing a segment family to the complete geometric graph on its
endpoints can produce an optimum whose restriction to the original
family is not optimal. The minimum contribution of the added edges can
also depend on the prescribed source colours.
\end{proposition}
\begin{proof}
Take $p_i=(i,i^2)$ for $0\leq i<6$ and abbreviate $p_ip_j$ by $ij$.
The source edges are $01,24,35$. Their only crossing pair is $24,35$,
so the source optimum is zero.

Every five-point subset is strictly convex. The crossing graph of its
five diagonals is an odd cycle, so every two-colouring contains a
monochromatic crossing. There are six five-point subsets, and every
crossing pair belongs to exactly two of them. Therefore the complete
six-point graph has at least three monochromatic crossing pairs in
every colouring.

Colour $03,04,13,14$ blue and all remaining edges red. Its monochromatic
crossing pairs are exactly $(02,15),(03,14),(24,35)$, so it is optimal
for the complete graph. Its restriction to the source has cost one.
In particular, when $24$ and $35$ have the same colour, the minimum
additional cost is two.

Alternatively colour $14,15,24,25$ blue and the other edges red.
The monochromatic crossing pairs are $(02,13),(04,35),(14,25)$.
This is again a complete optimum, but $24$ and $35$ now have different
colours. The source cost is zero and the minimum additional cost is
three. Hence the extra optimum is not independent of the source
colouring.
\end{proof}
